% ---
% titre: Lexique - Qui suis-je ? (fiches de tirage)
% theme: NO
% chapitre: Nombres naturels et décimaux
% sous_chapitre: 
% niveau: 9e	
% type: JEU
% tags: c-vocabulaire, p-entrainement
% difficulte: 1
% corrige: false
% source: latex
% ---
% NOTE : 24 fiches à imprimer et découper (3 mots par fiche : 2 mots du lexique NO
% + 1 mot du lexique GM). sous_chapitre laissé vide car le jeu porte sur le lexique
% cumulé de plusieurs chapitres. Tags à vérifier contre scripts/tags.md.
% ATTENTION : ne pas laisser de ligne vide entre deux \end{fiche} / \begin{fiche}
% à l'intérieur du tcbraster (la grille s'effondrerait en une seule colonne).

\newtcolorbox{fiche}[1]{%
  enhanced, sharp corners, colback=white, colframe=white, boxrule=0pt,
  height=2.9cm, left=3mm, right=3mm, top=1.5mm, bottom=1.5mm,
  before upper={\setlength{\parskip}{3pt}\textcolor{themeColor}{\textbf{\small Fiche #1}}\par},
  overlay={\draw[dashed, gray!70, line width=0.4pt] (frame.south west) rectangle (frame.north east);}
}

\header{Lexique --- Qui suis-je\,? Fiches}

\begin{tcbraster}[raster columns=3, raster column skip=0pt, raster row skip=0pt, raster valign=top]
\begin{fiche}{1}
\textbf{\large Multiple}\par
\textbf{\large Somme}\par
\textbf{\large Grandeur}
\end{fiche}
\begin{fiche}{2}
\textbf{\large Puissance}\par
\textbf{\large Ordre décroissant}\par
\textbf{\large Figure usuelle}
\end{fiche}
\begin{fiche}{3}
\textbf{\large Différence}\par
\textbf{\large Nombre}\par
\textbf{\large Aire}
\end{fiche}
\begin{fiche}{4}
\textbf{\large Associativité}\par
\textbf{\large Multiple commun}\par
\textbf{\large Base}
\end{fiche}
\begin{fiche}{5}
\textbf{\large Multiplication}\par
\textbf{\large Pgdc}\par
\textbf{\large Triangle}
\end{fiche}
\begin{fiche}{6}
\textbf{\large Chiffre}\par
\textbf{\large Terme}\par
\textbf{\large Rectangle}
\end{fiche}
\begin{fiche}{7}
\textbf{\large Nombre premier}\par
\textbf{\large Exposant}\par
\textbf{\large Sommet}
\end{fiche}
\begin{fiche}{8}
\textbf{\large Pgdc}\par
\textbf{\large Ordre décroissant}\par
\textbf{\large Parallèle}
\end{fiche}
\begin{fiche}{9}
\textbf{\large Somme}\par
\textbf{\large Division}\par
\textbf{\large Diagonale}
\end{fiche}
\begin{fiche}{10}
\textbf{\large Nombre naturel}\par
\textbf{\large Distributivité}\par
\textbf{\large Longueur}
\end{fiche}
\begin{fiche}{11}
\textbf{\large Nombre décimal}\par
\textbf{\large Dividende}\par
\textbf{\large Unité de mesure}
\end{fiche}
\begin{fiche}{12}
\textbf{\large Facteur}\par
\textbf{\large Commutativité}\par
\textbf{\large Préfixe}
\end{fiche}
\begin{fiche}{13}
\textbf{\large Chiffre}\par
\textbf{\large Quotient}\par
\textbf{\large Polygone}
\end{fiche}
\begin{fiche}{14}
\textbf{\large Nombre naturel}\par
\textbf{\large Diviseur commun}\par
\textbf{\large Parallélogramme}
\end{fiche}
\begin{fiche}{15}
\textbf{\large Élément neutre}\par
\textbf{\large Nombre premier}\par
\textbf{\large Figure complexe}
\end{fiche}
\begin{fiche}{16}
\textbf{\large Soustraction}\par
\textbf{\large Diviseur}\par
\textbf{\large Losange}
\end{fiche}
\begin{fiche}{17}
\textbf{\large Addition}\par
\textbf{\large Nombre}\par
\textbf{\large Quadrilatère}
\end{fiche}
\begin{fiche}{18}
\textbf{\large Terme}\par
\textbf{\large Arrondi}\par
\textbf{\large Hauteur}
\end{fiche}
\begin{fiche}{19}
\textbf{\large Multiple commun}\par
\textbf{\large Produit}\par
\textbf{\large Côté}
\end{fiche}
\begin{fiche}{20}
\textbf{\large Dividende}\par
\textbf{\large Multiplication}\par
\textbf{\large Périmètre}
\end{fiche}
\begin{fiche}{21}
\textbf{\large Nombre décimal}\par
\textbf{\large Ppmc}\par
\textbf{\large Carré}
\end{fiche}
\begin{fiche}{22}
\textbf{\large Facteur}\par
\textbf{\large Ppmc}\par
\textbf{\large Conversion d'unité}
\end{fiche}
\begin{fiche}{23}
\textbf{\large Division}\par
\textbf{\large Multiple}\par
\textbf{\large Perpendiculaire}
\end{fiche}
\begin{fiche}{24}
\textbf{\large Ordre croissant}\par
\textbf{\large Addition}\par
\textbf{\large Trapèze}
\end{fiche}
\end{tcbraster}